\documentclass[a4paper]{article}
% Español (México) con XeLaTeX: fontspec para fuentes Unicode, polyglossia
% para la división silábica y los nombres del documento en la variante mexicana.
\usepackage{fontspec}
\usepackage{polyglossia}
\setdefaultlanguage[variant=mexican]{spanish}
% Los nombres chinos van como «pinyin (original)»: xeCJK da a los caracteres
% Han su propia fuente; sin ella XeLaTeX los omite con un aviso.
% FandolSong no cubre algunos caracteres de nombres (U+73FA); WenQuanYi Zen Hei sí.
\usepackage[AutoFallBack=true]{xeCJK}
\setCJKmainfont{FandolSong-Regular.otf}
\setCJKfallbackfamilyfont{\CJKrmdefault}{WenQuanYi Zen Hei}
% polyglossia no traduce los nombres de listings.
\AtBeginDocument{\providecommand{\lstlistingname}{}\renewcommand{\lstlistingname}{Listado}%
  \providecommand{\lstlistlistingname}{}\renewcommand{\lstlistlistingname}{Índice de listados}}
\usepackage{amsmath,amssymb}
\usepackage{graphicx}
\usepackage[margin=2.35cm]{geometry}
\usepackage[most]{tcolorbox}
\usepackage{etoolbox}
\usepackage{listings}
\usepackage{xcolor}
\usepackage{hyperref}
\usepackage{booktabs,longtable,tabularx,array}
\usepackage{subcaption}
\usepackage{float}
\usepackage{enumitem}
\usepackage{microtype}
\usepackage{fancyhdr}
\usepackage{needspace}

\definecolor{kbBack}{HTML}{F2F6FA}
\definecolor{kbFrame}{HTML}{9DB4CC}
\definecolor{kbTitle}{HTML}{52708F}
\definecolor{ibBack}{HTML}{FBF5E7}
\definecolor{ibFrame}{HTML}{C9A961}
\definecolor{ibTitle}{HTML}{7D5F1E}
\definecolor{wbBack}{HTML}{FCF2F0}
\definecolor{wbFrame}{HTML}{D6A096}
\definecolor{wbTitle}{HTML}{9E4F43}
\definecolor{dbBack}{HTML}{F4F7F3}
\definecolor{dbFrame}{HTML}{AEC2AC}
\definecolor{dbTitle}{HTML}{5F7F64}

\tcbset{softbox/.style={
  enhanced,breakable,boxrule=0.8pt,arc=2pt,left=3mm,right=3mm,
  top=2mm,bottom=2mm,coltitle=white,fonttitle=\bfseries,
  toptitle=0.6mm,bottomtitle=0.6mm
}}
\newtcolorbox{knowledgebox}[1]{softbox,colback=kbBack,colframe=kbFrame,colbacktitle=kbTitle,title={#1}}
\newtcolorbox{importantbox}[1]{softbox,colback=ibBack,colframe=ibFrame,colbacktitle=ibTitle,title={#1}}
\newtcolorbox{warningbox}[1]{softbox,colback=wbBack,colframe=wbFrame,colbacktitle=wbTitle,title={#1}}
\newtcolorbox{dialoguebox}[1]{softbox,colback=dbBack,colframe=dbFrame,colbacktitle=dbTitle,title={#1}}

\lstset{
  language=Python,basicstyle=\ttfamily\small,
  keywordstyle=\color{kbTitle}\bfseries,stringstyle=\color{wbTitle},
  commentstyle=\color{dbTitle},breaklines=true,frame=single,numbers=left,
  numberstyle=\tiny\color{gray},captionpos=b,columns=fullflexible,
  keepspaces=true,showstringspaces=false,extendedchars=true
}
\BeforeBeginEnvironment{lstlisting}{\Needspace{12\baselineskip}}

\hypersetup{colorlinks=true,linkcolor=kbTitle,urlcolor=kbTitle,citecolor=wbTitle}
\setlist{itemsep=0.25em,topsep=0.4em}
\setlength{\parindent}{2em}
\setlength{\parskip}{0.25em}
\newcolumntype{L}[1]{>{\raggedright\arraybackslash}p{#1}}

\providecommand{\notetitle}{Notas del curso: [tema del curso]}
\providecommand{\noteauthors}{Elaboradas a partir de los materiales públicos de Stanford CS25}
\providecommand{\notedate}{Versión reescrita en 2026}
\providecommand{\videochannel}{Stanford Online}
\providecommand{\videopublishdate}{2022}
\providecommand{\videoduration}{[duración del video]}
\providecommand{\videourl}{}
\providecommand{\videocoverpath}{cover.jpg}
\providecommand{\coursesite}{https://web.stanford.edu/class/cs25/}
\providecommand{\repourl}{https://github.com/hqhq1025/ai-course-notes}
\providecommand{\skillversion}{wdkns/wdkns-skills@39f1a04}

\newcommand{\figsource}[1]{\par\vspace{-0.3em}{\footnotesize\color{gray}\emph{Fuente: #1}}\par}
\newcommand{\teachervoice}[1]{\begin{knowledgebox}{Nota de clase}#1\end{knowledgebox}}
\newcommand{\term}[1]{\textbf{#1}}

\pagestyle{fancy}
\fancyhf{}
\fancyfoot[C]{\thepage}
\fancyfoot[R]{\footnotesize\color{gray}\href{\repourl}{AI Course Notes}}
\renewcommand{\headrulewidth}{0pt}
\renewcommand{\footrulewidth}{0pt}

\newcommand{\makecscover}{
\begin{titlepage}
\centering
\vspace*{0.7cm}
{\Large Stanford CS25 · Transformers United\par}
\vspace{0.8cm}
{\huge\bfseries \notetitle\par}
\vspace{0.6cm}
{\large \noteauthors\par}
\vspace{0.25cm}
{\large \notedate\par}
\vspace{0.9cm}
\ifdefempty{\videocoverpath}{}{
  \includegraphics[width=0.86\textwidth,height=0.43\textheight,keepaspectratio]{\videocoverpath}\par
}
\vfill
\begin{tcolorbox}[width=0.92\textwidth,colback=black!2!white,colframe=black!45,boxrule=0.7pt,arc=2pt]
\textbf{Autor o canal del video}: \videochannel\par
\textbf{Fecha de publicación}: \videopublishdate\par
\textbf{Duración del video}: \videoduration\par
\textbf{Sitio del curso}: \href{\coursesite}{\nolinkurl{\coursesite}}\par
\ifdefempty{\videourl}{}{\textbf{Enlace del video}: \href{\videourl}{\nolinkurl{\videourl}}\par}
\textbf{Normas de elaboración}: \skillversion\ + las normas source-first / teacher-voice / visual-QA de este repositorio
\end{tcolorbox}
\end{titlepage}
\tableofcontents
\newpage
}

\usepackage{multicol}

\renewcommand{\notetitle}{Lecture 30: Near-Shallow Transformers y Precision Language Models}
\renewcommand{\noteauthors}{Elaborado a partir de la clase de Jake Williams, las slides recuperadas de la grabación oficial y subtítulos hechos manualmente}
\renewcommand{\notedate}{CS25 V4 · versión reescrita source-first de 2026}
\renewcommand{\videochannel}{Stanford Online}
\renewcommand{\videopublishdate}{Clase: 2024-05-02; publicación: 2024-05-23}
\renewcommand{\videoduration}{1:19:56}
\renewcommand{\videourl}{https://www.youtube.com/watch?v=zL9B3eXq0gY}
\renewcommand{\videocoverpath}{cover.jpg}

\newcommand{\lecturefigure}[3]{%
\begin{figure}[H]
  \centering
  \includegraphics[width=0.97\textwidth,height=0.49\textheight,keepaspectratio]{slides-images/#1}
  \caption{#2}
  \figsource{#3}
\end{figure}
}

\let\standardtableofcontents\tableofcontents
\renewcommand{\tableofcontents}{%
  \begingroup
  \small
  \setlength{\parskip}{0pt}
  \standardtableofcontents
  \endgroup
}

\begin{document}
\makecscover

\section{Ubicación de la clase: no se trata simplemente de «hacer menos profundo el Transformer»}

El término «casi superficial» (near-shallow) del título puede hacer creer que la clase solo trata de reducir el número de capas del Transformer. En realidad, el ponente propone una ruta más concreta y más sujeta a condiciones. Primero, reescribe una parte del trabajo de «proyectar primero y comparar después» de la self-attention estándar como «transformar directamente las comparaciones entre vectores»; segundo, usa la generalized co-occurrence (coocurrencia generalizada) para construir un warm start no aleatorio para la softmax layer; tercero, cuando el embedding está fijo, almacena en caché las comparaciones entre vectores de la capa inferior, de modo que una parte del cálculo ya no se repite en cada entrenamiento o inferencia; cuarto, observa si esta estructura permite un tiny controller entrenado solo con datos locales de la aplicación. Aquí, near-shallow significa que \textbf{el cálculo reutilizable de la capa inferior hace menos profunda la ruta efectiva en línea}, no que a la red le quede una sola capa.

\subsection{Primero se corrigen las fuentes: el título y el contenido de las notas anteriores no eran confiables}

La versión anterior fechaba la clase en 2026, usaba una URL de una Stanford lecture inexistente y añadía mucho contenido como drift gate, attention observability board, fairness dashboard, weekly newsletter y RLHF governance roadmap. Estos mecanismos no aparecen en la página oficial del curso ni en la grabación, ni tampoco en los artículos relacionados del ponente. Esta clase toma ahora como única base la sesión del 2 de mayo de 2024, la grabación oficial de Stanford Online \texttt{zL9B3eXq0gY}, los subtítulos oficiales \texttt{en-US} hechos por personas y los 27 slide states didácticos recuperados de la grabación.

\lecturefigure{slide-01-motivation.jpg}{Pregunta central: ¿se puede aprender directamente la transformación de «resultados de comparación a feature weights»?}{Grabación oficial de Stanford Online 00:04:39}

La primera slide didáctica compara dos rutas. La attention estándar cambia el espacio de representación mediante $W_q,W_k$ y luego produce el score con el dot product; la forma alternativa del ponente conserva la comparación entre los vectores originales $XX^\top$ y usa otra matriz $W$ para convertir esas comparaciones en attention logits. Toda la inicialización explícita, el almacenamiento en caché y los experimentos de edge que siguen giran en torno a esta reescritura.

\begin{importantbox}{Las tres preguntas centrales de esta clase}
\begin{enumerate}
  \item \textbf{Initialization:} ¿se puede evitar arrancar con parámetros aleatorios y obtener un valor inicial útil directamente a partir de estadísticas de entrada y salida?
  \item \textbf{Reusable computation:} ¿qué representaciones o comparaciones entre vectores pueden congelarse y almacenarse en caché para reducir el cálculo repetido?
  \item \textbf{Application-only learning:} ¿puede un modelo pequeño aprender solo de los datos de interacción de una tarea local, sin depender de un preentrenamiento general?
\end{enumerate}
\end{importantbox}

\teachervoice{00:01:19--00:02:49, Williams explica primero su formación en matemáticas, física y quantitative linguistics. Parte deliberadamente de la estructura estadística y de fórmulas explícitas, y no de las scaling recipes habituales.}

\subsection{Niveles de evidencia: el conocimiento estándar, las derivaciones de los artículos y los experimentos de clase no deben mezclarse}

En esta clase aparecen a la vez fórmulas estándar del Transformer, derivaciones de preprints del equipo, resultados experimentales en las slides de clase, una demostración de viabilidad en Le Potato y planteamientos a futuro en la sesión de Q\&A. Para no presentar un prototipo de investigación como una receta consolidada, conviene mantener la siguiente separación por niveles al leer.

\begin{knowledgebox}{Registro de evidencia}
\begin{tabularx}{\textwidth}{L{0.17\textwidth} L{0.29\textwidth} X}
\toprule
Nivel & Ejemplos & Lectura correcta \\
\midrule
Antecedentes estándar & softmax, dot-product attention, perplejidad, padding & Puede tomarse como base técnica común \\
Derivaciones de artículos & solución explícita de generalized co-occurrence, SAFFU hidden targets & Deben respetarse los supuestos y las condiciones de aproximación del artículo \\
Experimentos de clase & curvas warm/cold, MNIST, BabyLM, switch task & Solo respaldan lo observado en la configuración correspondiente \\
Juicios del ponente & el packing no es el context model correcto, nombre de PLM & Son posturas de investigación, no consenso \\
Trabajo futuro & convolution warm start, multimodal, RLHF, implementación pública & Aún no está terminado; no debe presentarse como capacidad verificada \\
\bottomrule
\end{tabularx}
\end{knowledgebox}

\begin{warningbox}{SAFFU no es una arquitectura estándar disponible}
Self-Attentive Feed-Forward Unit (SAFFU, unidad de prealimentación autoatenta) y Precision Language Model (PLM, modelo de lenguaje de precisión) son términos de investigación que usa el equipo del ponente. No son módulos estandarizados en PyTorch/Hugging Face, y tampoco puede concluirse, solo con base en esta clase, que «entrenar un LLM ya no requiere backpropagation».
\end{warningbox}

\subsection{Resumen de la sección}

La clase real gira en torno a la inicialización no aleatoria, el almacenamiento en caché y los modelos locales pequeños, no en torno a los procesos de gobernanza de la versión anterior. Estas notas ubicarán cada resultado dentro de sus supuestos concretos: softmax, características no negativas, embedding fijo, datos limitados y un prototipo de edge en etapa temprana.
\section{De la attention estándar a SAFFU: qué paso se modifica exactamente}

Para entender el mecanismo alternativo, primero hay que descomponer la self-attention estándar en cuatro pasos: «proyección de la representación», «similitud», «normalización» y «agregación ponderada». El ponente no niega el valor de la Q/K projection estándar; lo que plantea es: cuando las entradas ya se encuentran en el mismo espacio de baja dimensión, ¿puede aplicarse directamente una transformación aprendible a la propia matriz de comparación, para reducir las proyecciones adicionales y los parámetros aleatorios?

\subsection{Forma estándar: primero proyectar, después comparar}

Sean los hidden states de una secuencia $X\in\mathbb{R}^{L\times d}$; la attention estándar de una sola cabeza se escribe como
\[
Q=XW_q,\qquad K=XW_k,\qquad V=XW_v,
\]
\[
A_{\text{std}}=\operatorname{softmax}\!\left(\frac{QK^\top}{\sqrt{d_a}}+M\right),
\qquad H=A_{\text{std}}V.
\]
Aquí $L$ es el número de tokens, $d$ es la dimensión de entrada, $d_a$ es la attention dimension y $M$ es la máscara causal o de padding. La tarea de $W_q,W_k$ no es simplemente «calcular la similitud», sino transformar primero la entrada a un espacio de comparación más adecuado para la tarea de predicción en cuestión.

\begin{knowledgebox}{Por qué la proyección suele ser útil}
Cuando query y key provienen de modalidades distintas, de dimensiones distintas o de roles semánticos distintos, la cross-attention debe llevarlas primero a un espacio comparable. Incluso en la self-attention, la proyección permite que distintas heads aprendan similitudes bilineales distintas. La ganancia de eficiencia del mecanismo alternativo proviene de reducir este paso, pero al mismo tiempo renuncia a una parte de la libertad de representación.
\end{knowledgebox}

\teachervoice{00:02:54--00:05:06, el ponente describe la attention estándar como re-dimensionalization: Q/K cambian el espacio de representación para que el dot product pueda producir pesos útiles para la tarea; la forma alternativa, en cambio, aprende directamente cómo interpretar la comparación entre los vectores existentes.}

\subsection{Definición de SAFFU: transformar directamente la comparison matrix}

SAFFU supone que los vectores de entrada ya se encuentran en un mismo espacio comparable. Primero se calcula la matriz de comparación original $C=XX^\top$ y luego se transforma la comparación con una matriz aprendible $W$, con lo que se obtiene
\[
A_{\text{saffu}}=\operatorname{softmax}(XWX^\top+M),
\qquad H=A_{\text{saffu}}X.
\]
Aquí $W\in\mathbb{R}^{d\times d}$ define directamente una bilinear form (forma bilineal). Si $W=I$, el modelo se reduce al dot product en el espacio original; en general, $W$ puede rotar, escalar o recombinar las interacciones entre dimensiones, pero no proyecta query y key por separado en dos espacios independientes.

\lecturefigure{slide-02-saffu-definition.jpg}{Definición, supuestos y restricciones de representación de SAFFU}{Video oficial de Stanford Online 00:06:00}

\begin{importantbox}{Relación algebraica entre la attention estándar y SAFFU}
Puesto que
\[
QK^\top=XW_qW_k^\top X^\top,
\]
si se denota $W=W_qW_k^\top$, ambas formas pueden ponerse en correspondencia en la forma bilineal del score. La diferencia está en la parametrización: la forma estándar descompone $W$ en dos projections, mientras que SAFFU trabaja directamente con la matriz combinada y deriva a partir de ella una inicialización explícita. La parametrización directa no es automáticamente más potente ni automáticamente más estable.
\end{importantbox}

\begin{warningbox}{«Sin Q/K» no equivale a sin transformación matricial}
SAFFU sigue teniendo $W$ y sigue aprendiendo cómo debe modificarse la comparación entre tokens. Lo que elimina es la parametrización con dos projections independientes y la representación intermedia, no convierte la attention en un promedio sin parámetros.
\end{warningbox}

\teachervoice{00:05:15--00:05:43, Williams afirma explícitamente que el mecanismo alternativo y la projection tradicional pueden usarse al mismo tiempo; en esta clase se retira temporalmente la projection para aislar y estudiar el propio mecanismo comparison-to-weight.}

\subsection{General design: de la comparación y la agregación a la salida}

El diagrama de la arquitectura usa $W$ para la attention comparison transform, $U$ para la transformación decoder/feed-forward aplicada al hidden state agregado y $O$ para el mapeo de salida de la capa superior. Al leer la figura, no hay que llamar a los tres «parámetros de attention»: $W$ determina los pesos de las posiciones, $U$ modifica las features agregadas y $O$ mapea el hidden state al objetivo de predicción.

\lecturefigure{slide-03-general-saffu-design.jpg}{General SAFFU encoder-decoder: reparto de funciones entre $W$, $U$ y el mapeo de salida}{Video oficial de Stanford Online 00:07:38}

Para la $m$-ésima muestra, sea la context matrix $X_m\in\mathbb{R}^{K\times d}$ y elíjase la fila query $x_{m,h}$; entonces puede escribirse
\[
a_m=\operatorname{softmax}(x_{m,h}WX_m^\top),
\]
\[
h_m=a_mX_mU,\qquad \hat y_m=\operatorname{softmax}(h_mO).
\]
$K$ es el número de features/tokens en el context, $a_m\in\mathbb{R}^{K}$ es la attention distribution, y $U$ y $O$ se encargan, respectivamente, de la transformación oculta y de la salida final de clasificación o de modelado de lenguaje.

\begin{knowledgebox}{Qué significa exactamente near-shallow}
El modelo puede seguir teniendo varias capas. Near-shallow significa que algunas representaciones y pairwise comparisons de las capas inferiores pueden determinarse de antemano, almacenarse en caché o inicializarse explícitamente de forma local, de modo que el aprendizaje en línea ocurra sobre todo en un número menor de parámetros de las capas superiores. Describe las dependencias de cómputo y las rutas de actualización, no simplemente un número de depth.
\end{knowledgebox}

\teachervoice{00:05:43--00:07:10, el ponente relaciona la reducción de projections y de parámetros aleatorios con near-shallow, pero reconoce al mismo tiempo que este diseño depende más de que los vectores de entrada tengan la misma dimensión y estén bien organizados; no es un reemplazo incondicional.}

\subsection{Resumen de la sección}

La modificación central de SAFFU es parametrizar directamente $XWX^\top$ e incorporar esta forma a una inicialización explícita. La projection de la attention estándar es más flexible; el atractivo de SAFFU proviene de tener menos parámetros aleatorios independientes, valores iniciales que pueden derivarse y comparaciones que potencialmente pueden almacenarse en caché. Las ganancias que se presentan más adelante deben leerse junto con estas restricciones.
\section{Inicialización no aleatoria: por qué el bit-cipher aparece antes que la attention}

La inicialización explícita no puede derivarlo todo a partir de embeddings aleatorios. Si los vectores de entrada son aleatorios, tienen signos no controlados o cambian en cada actualización, las estadísticas de coocurrencia y las cachés posteriores pierden una referencia estable. Por eso el curso analiza primero la intuición de Lottery Ticket, la necesidad de representaciones de baja dimensión y el bit-cipher, y después pasa al warm start de la attention.

\subsection{Lottery Ticket es una motivación, no una demostración en contra de multi-head}

Varias attention heads tienen la misma estructura y el mismo objetivo de entrenamiento, pero la inicialización aleatoria hace que adopten papeles distintos durante la optimización. El ponente vincula esta incertidumbre con la Lottery Ticket Hypothesis: una red grande puede contener unas cuantas subestructuras realmente eficaces, mientras que el costo de entrenar una gran cantidad de parámetros aleatorios no termina por convertirse en capacidad aprovechable.

\lecturefigure{slide-04-no-lottery-tickets.jpg}{Transforming without lottery tickets: reducir la inicialización aleatoria y las copias simétricas}{Grabación oficial de Stanford Online 00:11:04}

\begin{warningbox}{No se puede concluir directamente de Lottery Ticket que «las múltiples cabezas son inútiles»}
Las distintas heads pueden aprender relaciones complementarias y también pueden romper la simetría mediante el entrenamiento. La slide ofrece una motivación de investigación: si fuera posible encontrar un buen punto de partida directamente a partir de las estadísticas de los datos, ¿se podría depender menos de un sorteo aleatorio? No es una conclusión de un estudio de ablación sobre multi-head attention.
\end{warningbox}

\teachervoice{00:08:03--00:10:44, Williams explica la diferenciación aleatoria de papeles entre varias heads de la misma estructura como un problema de tipo lottery ticket, y enfatiza que la inicialización aleatoria puede hacer que parámetros costosos terminen sin desempeñar una función estable.}

\subsection{La dimensionality reduction sigue siendo inevitable}

El problema de la subsección anterior es cómo reducir el «sorteo» aleatorio entre heads de la misma estructura, pero antes de la attention existe una fuente de aleatoriedad aún más temprana: el embedding de los tokens. Si no se estabiliza primero el sistema de coordenadas de la entrada, tanto la generalized co-occurrence como la comparison cache posteriores se desplazarán a medida que se actualice el embedding. Por eso esta subsección explica primero por qué es necesario comprimir la dimensión y después distingue entre «baja dimensión» y «una buena representación no aleatoria de baja dimensión». Sea $V$ el tamaño del vocabulario; si un token se representa directamente como one-hot $y\in\{0,1\}^{V}$, la dimensión de las matrices de pesos posteriores crece con $V$. La práctica habitual de aprender un embedding $E\in\mathbb{R}^{V\times d}$ toma $d\ll V$, pero un embedding aleatorio está muy lejos de la loss de salida, y el gradiente debe atravesar toda la red para que adquiera forma poco a poco.

\lecturefigure{slide-05-dimensionality-reduction.jpg}{Las representaciones de baja dimensión siguen siendo necesarias, pero un embedding aleatorio es difícil de estabilizar y de interpretar}{Grabación oficial de Stanford Online 00:11:21}

\begin{importantbox}{La compresión de la representación y la inicialización son dos problemas distintos}
\textbf{Dimensionality reduction} determina cuántas dimensiones se usan para expresar un token; \textbf{initialization} determina qué representan esas dimensiones desde el principio. Reducir la dimensión de $V$ a $d$ no produce automáticamente una buena representación: una matriz aleatoria solo le da al optimizador un punto de partida entrenable.
\end{importantbox}

\teachervoice{00:11:14--00:12:40, el ponente enfatiza que, con vocabularios grandes, la representación de baja dimensión es una necesidad computacional, pero el embedding es lo más alejado de la loss de salida, y un punto de partida aleatorio requiere un entrenamiento costoso para estabilizarse.}

\subsection{Bit-cipher: primero garantizar la unicidad, después hablar de semántica}

El bit-cipher codifica los tokens con un bit vector de longitud $b$. Excluyendo el vector de puros ceros, en teoría hay como máximo
\[
N_{\text{codes}}=2^b-1
\]
codificaciones únicas, por lo que cubrir el vocabulario requiere al menos $b\ge \lceil\log_2(V+1)\rceil$. El equipo asigna los bit patterns ordenando los tokens por unigram frequency, de modo que los tokens de alta frecuencia reciben primero codificaciones dispersas y fáciles de distinguir, y después se normaliza
\[
e_i=\frac{\eta_i}{\|\eta_i\|_1},\qquad \eta_i\in\{0,1\}^{b}.
\]

\lecturefigure{slide-06-bit-cipher.jpg}{Bit-cipher: ampliar el one-hot encoding con bit vectors ordenados por frecuencia}{Grabación oficial de Stanford Online 00:15:01}

\begin{knowledgebox}{Qué resuelve el bit-cipher y qué no resuelve}
Proporciona una token identity determinista, única, de baja dimensión y no negativa, lo que facilita las fórmulas explícitas y la reproducción de experimentos. Por sí mismo no garantiza que los vectores de palabras semánticamente cercanas estén próximos; después todavía hay que recurrir a la agregación de co-occurrence o a las radial context statistics para añadir estructura de relaciones.
\end{knowledgebox}

\begin{warningbox}{Ordenar por frecuencia no es un teorema semántico}
Que los tokens de alta frecuencia reciban codificaciones más simples es una decisión de diseño de ingeniería. La frecuencia por sí sola no permite inferir similitud semántica, papel sintáctico ni alineación entre idiomas; eso debe verificarse mediante tareas posteriores y análisis de correlación.
\end{warningbox}

\teachervoice{00:12:53--00:15:45, el ponente explica que el verdadero valor del bit-cipher es la deterministic low dimensionalization. Parece solo un paso intermedio, pero es el punto de partida real que permitió al equipo seguir derivando la inicialización no aleatoria de SAFFU.}

\subsection{Términos clave: cinco palabras que se repetirán más adelante}

Estos términos suelen condensarse en una sola frase, «closed-form training», pero representan etapas distintas y afirmaciones de distinto alcance.

\begin{knowledgebox}{Glosario de inicialización y optimización}
\begin{tabularx}{\textwidth}{L{0.18\textwidth} L{0.27\textwidth} X}
\toprule
Término & Definición & Papel en esta clase \\
\midrule
Cold start & Los parámetros parten de una distribución aleatoria & Sirve como baseline del entrenamiento convencional \\
Warm start & Los parámetros parten de valores no aleatorios obtenidos de estadísticas de los datos & Reduce la loss/perplejidad inicial y mejora la trayectoria posterior \\
Explicit solution & Parámetros aproximados calculados directamente a partir de estadísticos, sin gradientes iterativos & Inicializa la softmax layer y los componentes de SAFFU \\
Priming number & Cantidad de normalización que, en la solución explícita, describe la norma o el número de features de la entrada & Ajusta la column translation de la matriz de coocurrencia \\
Label embedding & Representación de baja dimensión del objetivo de predicción & Proporciona hidden targets calculables para las capas intermedias \\
\bottomrule
\end{tabularx}
\end{knowledgebox}

\subsection{Resumen de la sección}

La inicialización no aleatoria requiere entradas estables, no negativas y de baja dimensión. El bit-cipher resuelve primero una identity encoding reproducible y así le da a la generalized co-occurrence un sistema de coordenadas fijo. No es el punto final de la representación semántica, sino una de las condiciones previas del warm start explícito de SAFFU.
\section{Warm start explícito de softmax: de la matriz de coocurrencias a los parámetros}

Esta sección es el núcleo matemático de la clase. El objetivo no es afirmar que toda neural network tiene una solución exacta en forma cerrada, sino mostrar que, para non-negative features y una single layer con activación softmax, se puede obtener un buen valor inicial aproximado a partir de los outer products de entrada y salida; después, esa idea se extiende a varios componentes de SAFFU.

\subsection{La generalized co-occurrence no se limita a conteos en ventanas de palabras}

Sea $H\in\mathbb{R}_{\ge 0}^{M\times D}$ la feature matrix de entrada de $M$ muestras, y sea el objetivo una distribución one-hot o de probabilidad $Y\in\mathbb{R}_{\ge 0}^{M\times N}$. La coocurrencia generalizada se define como
\[
F(H,Y)=\sum_{m=1}^{M}H_{m,:}\otimes Y_{m,:}=H^\top Y,
\]
donde $F_{j,i}$ acumula la intensidad con la que la $j$-ésima feature de entrada y el $i$-ésimo target de salida aparecen juntos. Puede provenir del token context, de los píxeles de una imagen, de las salidas de una capa oculta o de cualquier otra feature no negativa; no tiene que ser una tabla de frecuencias de palabras con una ventana de radio tradicional.

\lecturefigure{slide-07-softmax-cooccurrences.jpg}{Softmax y generalized co-occurrence: fórmula de inicialización explícita}{Grabación oficial de Stanford Online 00:19:13}

Para un softmax decoder $U\in\mathbb{R}^{D\times N}$, la forma aproximada que presenta la slide puede escribirse como
\[
U_{j,i}\approx \log\!\bigl(F(H,Y)_{j,i}+\epsilon\bigr)
-\frac{K-1}{K}\log\!\left(\sum_{d=1}^{D}F(H,Y)_{d,i}+\epsilon\right),
\]
donde $K$ es el priming number y $\epsilon>0$ evita tomar numéricamente el logarithm de cero. El primer término conserva la asociación feature--target; el segundo normaliza por columna de salida los target de alta frecuencia.

\begin{importantbox}{Intuición de la fórmula: el log invierte aproximadamente el softmax}
El softmax usa exponentials para convertir los logits en probabilidades; la solución explícita toma el log de la intensidad de coocurrencia, lo que equivale a regresar de la escala de probabilidad a la escala de logit. La column translation no cambia la estructura relativa del softmax dentro de una misma columna de salida, pero sí compensa la frecuencia global de los distintos target.
\end{importantbox}

\begin{warningbox}{Esta no es la solución óptima exacta para cualquier layer}
Los supuestos principales son: features no negativas, activation softmax, estadísticos que puedan estimarse de forma estable y una norm de entrada que pueda describirse con $K$. La composition de varias capas, ReLU/GELU, las negative features y los attention hidden targets requieren derivaciones adicionales.
\end{warningbox}

\teachervoice{00:15:55--00:18:09: el ponente subraya que aquí la co-occurrence es la outer-product sum de inputs y outputs arbitrarios, no limitada a las ventanas tradicionales de NLP; el término relacionado con $K$ expresa de forma aproximada la corrección que el número de context/feature introduce en la probabilidad condicional.}

\subsection{Priming number: aproximación del número de features y de la norm}

Si cada vector de entrada tiene exactamente $K$-norm, es decir, $\|H_{m,:}\|_1=K$, $K$ puede entrar directamente en la normalización. En la práctica, las features continuas rara vez cumplen una norm uniforme; la clase y el artículo usan el promedio
\[
\widehat K=\frac{1}{M}\sum_{m=1}^{M}\|H_{m,:}\|_1
\]
como aproximación. Esta estimación muestra que la solución explícita sigue necesitando una data-dependent calibration; no es una inicialización universal totalmente libre de hiperparámetros.

\begin{knowledgebox}{Cómo verificar la aproximación del priming number}
Como mínimo, hay que reportar la distribución de $\|H_m\|_1$, su media y varianza, la proporción de valores cero, el lugar donde se añade $\epsilon$ y la sensibilidad de la validation loss a distintos valores de $K$. Con un solo promedio no es posible saber si una heavy-tail o los outlier rompen la aproximación.
\end{knowledgebox}

\teachervoice{00:31:20--00:32:15: Williams explica que el priming number para features no unitarias puede estimarse con la norm promedio de la entrada, y reconoce que a la slide le falta la figura correspondiente; la evidencia completa debe verificarse en el artículo.}

\subsection{De una sola capa a SAFFU: la dificultad está en los hidden targets}

En el feed-forward decoder, la entrada $H$ y el target $Y$ están definidos explícitamente, así que $F(H,Y)$ se puede calcular directamente. La attention, en cambio, no tiene etiquetas explícitas que le indiquen al modelo «a qué posición debe prestar atención». La extensión que propone el equipo consiste en derivar primero un hidden target: hacer que la representación que produce la attention se acerque lo más posible a la entrada que espera la capa superior $U$, y luego inicializar capa por capa, de abajo hacia arriba.

Si se simplifica la notación como
\[
X\xrightarrow{W} A\xrightarrow{} Z=AX
\xrightarrow{U} H\xrightarrow{O}\hat Y,
\]
el orden del warm start es: estabilizar $X$; estimar los attention targets a partir de lo que requiere la capa superior e inicializar $W$; inicializar $U$ con $Z$ y los label embeddings; por último, inicializar $O$ con $H$ y el $Y$ real.

\begin{warningbox}{Un local fit no equivale a un compositional optimum}
La inicialización explícita capa por capa solo aprovecha inputs/targets locales. La backpropagation, mediante la chain rule, propaga los efectos de orden superior de la tarea final sobre todas las capas; por eso, después del warm start todavía puede hacer falta optimización iterativa. Lo que sostiene la clase es «un mejor punto de partida», no que «la composition se haya resuelto en forma cerrada».
\end{warningbox}

\teachervoice{00:18:13--00:20:57: el ponente define la dificultad de la attention como la ausencia de un target vector directo; el equipo construye los hidden targets analizando lo que esperan las capas superiores y, con ayuda de los bit-cipher label embeddings, reduce los objetivos intermedios a una dimensión manejable.}

\subsection{Procedimiento ejecutable: acumular estadísticos a partir de parámetros en cero}

La slide 8 condensa la teoría en un warm-start procedure. En términos de ingeniería, lo más importante no es que la fórmula parezca «closed form», sino que solo requiere una forward-style accumulation: recorrer los datos, acumular outer products y luego aplicar la normalización y la log transform.

\lecturefigure{slide-08-nonrandom-feedforward-init.jpg}{Pasos de cálculo de la non-random feed-forward initialization}{Grabación oficial de Stanford Online 00:21:16}

\begin{lstlisting}[caption={Pseudocódigo shape-first del warm start explícito de una softmax layer}]
F = zeros(feature_dim, target_dim)
for features, targets in dataset:
    assert all(features >= 0) and all(targets >= 0)
    F += features.T @ targets

K_hat = mean_l1_norm(dataset.features)
column_mass = F.sum(axis=0, keepdims=True)
weights = log(F + eps) - ((K_hat - 1) / K_hat) * log(column_mass + eps)
return weights
\end{lstlisting}

\begin{knowledgebox}{Por qué es fácil de paralelizar}
Cada worker puede acumular de forma independiente su $F_r=H_r^\top Y_r$ local y, al final, se suma $F=\sum_rF_r$. Esto se parece a la agregación de sufficient statistics y no exige sincronizar la gradient trajectory completa después de cada mini-batch. Aun así, la memoria, el rango numérico y el almacenamiento disperso requieren diseño de ingeniería.
\end{knowledgebox}

\teachervoice{00:21:15--00:22:43: Williams recurre a «todos los parámetros parten de zero» para subrayar que el método no es aleatorio: se recorren los datos acumulando los outer products de entrada y salida, y después la normalization y el logarithm dan el valor inicial.}

\subsection{Resumen de la sección}

La conclusión reproducible más sólida del warm start explícito es la siguiente: bajo las condiciones específicas de softmax y features no negativas, se puede calcular una data-informed initialization mediante la generalized log-co-occurrence. Extenderla a SAFFU requiere hidden targets y composition capa por capa; la fórmula de una sola capa no puede generalizarse sin condiciones a todos los módulos de un Transformer moderno.
\section{Evidencia de warm start: cómo deben leerse las curvas}

El curso usa dos grupos de experimentos para mostrar que la inicialización puede ser útil entre tareas: las trayectorias de perplejidad de un modelo de lenguaje de una sola capa y la clasificación de MNIST con una y dos capas. Lo importante de las gráficas es el punto de partida y la trayectoria de optimización, no anunciar que la nueva arquitectura ya alcanzó la calidad de los modelos grandes.

\subsection{Modelo de lenguaje de una sola capa: el punto de partida es más bajo y la trayectoria sigue mejorando}

La sección anterior solo presentó el método para calcular los parámetros; esta sección pasa a la pregunta empírica más directa: ¿este valor inicial data-informed está realmente más cerca de un modelo utilizable que un punto de partida aleatorio? Al leer las curvas, primero se compara el punto de partida en la epoch 0 y después la trayectoria de descenso y la posición de early-stopping bajo la misma update rule; no se debe mirar primero solo el último punto e ignorar que las warm-start statistics también tienen un costo de cómputo.

La Perplexity (perplejidad) es la exponencial de la cross-entropy promedio:
\[
\operatorname{PPL}=\exp\!\left(-\frac{1}{T}\sum_{t=1}^{T}\log p(x_t\mid x_{<t})\right).
\]
Puede entenderse intuitivamente como el «número efectivo de candidatos» que el modelo enfrenta en cada token; mientras más baja, por lo general mejor. Sin embargo, la PPL obtenida con distinta tokenización, distintos datos o distinta configuración de context no puede compararse directamente.

\lecturefigure{slide-09-single-layer-warm-starts.jpg}{Trayectorias de entrenamiento con warm start y cold start en el modelo de una sola capa}{Grabación oficial de Stanford Online 00:24:03}

Al leer la figura, primero se observa la epoch 0: la PPL del cold start está cerca del vocabulary size, mientras que la del warm start ya es claramente menor. Después se observan la pendiente y el punto final: bajo la misma regla de iterative update, la curva de warm start sigue descendiendo y, al detenerse de forma temprana, mantiene una PPL más baja. Esto respalda que «un valor inicial explícito mejora el punto de partida y la optimización posterior», pero no controla todos los componentes de un Transformer moderno.

\begin{knowledgebox}{Lectura de la figura: esta gráfica no tiene self-attention}
En la sesión de Q\&A se aclaró que se trata de un feed-forward experiment con context one-hot concatenado. Las features de entrada ya están separadas, por lo que casi no se necesita attention para deshacer la superposition. Lo que valida es la single-layer explicit initialization, no una ablación completa de la attention de SAFFU.
\end{knowledgebox}

\begin{warningbox}{Usar la misma learning rate no equivale a una comparación completamente justa}
Warm start y random start pueden requerir distintos learning-rate schedule, regularization o paciencia de early-stopping. Una comparación rigurosa debe reportar al mismo tiempo una tuned cold baseline, varias seeds, el wall-clock y el costo de leer los datos y de calcular los estadísticos explícitos.
\end{warningbox}

\teachervoice{00:22:47--00:24:24, el ponente enfatiza que el warm start reduce la PPL desde antes del entrenamiento y que después sigue mejorando con la misma learning rate; si se usa early stopping, este también se activa antes y con una PPL más baja.}

\teachervoice{01:07:53--01:09:37, en la sesión de Q\&A se corrigió una lectura equivocada frecuente: esta gráfica warm/cold no tiene attention; el trabajo inicial era solo un feed-forward de una sola capa, y la solución explícita para la attention de SAFFU se agregó después.}

\subsection{MNIST: muestra que la fórmula a nivel de layer no se limita al texto}

Los pixels de MNIST son no negativos, así que el mismo warm start de la softmax layer puede aplicarse a image classification. En la figura, los modelos de una y dos capas con warm start alcanzan antes una accuracy alta; sin embargo, el modelo no incluye la convolution, la normalization ni el pipeline de augmentation de las redes de visión modernas.

\lecturefigure{slide-10-image-classification-warm-starts.jpg}{Curvas de clasificación con warm start en MNIST}{Grabación oficial de Stanford Online 00:30:45}

\begin{importantbox}{La conclusión mínima que respaldan los resultados de MNIST}
La softmax initialization explícita puede recibir features no lingüísticas y no negativas, y mejora el punto de partida en esta configuración de clasificación. No demuestra que la misma fórmula ya cubra convolution, ViT, diffusion ni datos visuales del mundo real.
\end{importantbox}

\teachervoice{00:29:14--00:31:20, el ponente usa MNIST para mostrar la layer-level generality de la fórmula; no afirma haber construido una vision architecture completa de alto rendimiento.}

\teachervoice{01:09:42--01:11:46, en la sesión de Q\&A se acotó todavía más el alcance: para manejar convolution u otras activation, todavía hay que derivar por separado el warm start de esas layer; de lo contrario, los parámetros aleatorios vuelven a introducir ruido.}

\subsection{Qué baselines deben agregarse al experimento}

En este tipo de investigación sobre initialization no basta con mirar la accuracy/PPL final. Hay que separar al menos cuatro tipos de costo: el cálculo de los estadísticos explícitos, el tuning de la baseline aleatoria, la backpropagation posterior y el efecto de fijar o actualizar los embedding sobre la caché.

\begin{knowledgebox}{Matriz experimental mínima}
\begin{tabularx}{\textwidth}{L{0.18\textwidth} L{0.24\textwidth} X}
\toprule
Comparación & Elementos fijos & Qué se debe reportar \\
\midrule
Cold vs warm & architecture, data, optimizer & seeds, loss inicial/final, wall-clock \\
Warm only vs warm+tune & initialization statistics & costo adicional de gradiente y ganancia en generalización \\
Frozen vs thawed embeddings & warm-start weights & aciertos de cache, velocidad de entrenamiento, calidad final \\
Local explicit vs end-to-end & parameter count & compositional gap y failure cases \\
\bottomrule
\end{tabularx}
\end{knowledgebox}

\subsection{Resumen de la sección}

Los experimentos de warm start muestran un mejor punto de partida y una mejor trayectoria de optimización, y en MNIST exhiben una capacidad limitada de transferirse a otro dominio. Lo que validan son condiciones específicas de layer y de datos, no que «todo el aprendizaje profundo pueda sustituirse por matrices de coocurrencia».
\section{Context, multi-context y near-shallow caching}

La primera mitad de la clase responde «cómo arrancan los parámetros»; esta sección responde «cómo repetir menos cómputo». El curso sostiene que hay un desajuste entre un context fijo muy largo, el document packing y los datos de longitud variable; la forma matricial de SAFFU permite elegir subbloques de parámetros según la context length y, combinada con un frozen embedding, almacenar en caché las pairwise comparisons de la capa inferior.

\subsection{Una ventana larga no contiene automáticamente más información útil}

Con una context window fija $K_{\max}$, el attention cost es aproximadamente
\[
C_{\text{attn}}=\Theta(BK_{\max}^{2}d),
\]
donde $B$ es el tamaño de lote y $d$ es la hidden dimension. Si la mayoría de los documentos tienen longitud $K_m\ll K_{\max}$, los padding tokens siguen ocupando posiciones en la matriz. Una ventana más larga aumenta la capacidad, pero no garantiza que cada ejemplo aporte realmente más context relevante.

\lecturefigure{slide-11-longer-contexts.jpg}{Una ventana larga es un límite de ingeniería arbitrario; no significa que cada ejemplo necesite el mismo context}{Grabación oficial de Stanford Online 00:34:22}

\begin{knowledgebox}{Context tiene al menos tres significados}
\textbf{Block context} es el número fijo de tokens que el modelo recibe de una vez; \textbf{document context} es el texto anterior y posterior realmente relevante dentro del mismo documento; \textbf{linguistic/radial context} es la vecindad estadística alrededor del token objetivo. Llamar context a los tres oculta las diferencias entre padding, packing y los supuestos de modelado.
\end{knowledgebox}

\teachervoice{00:32:20--00:34:55, Williams considera que una ventana más larga aumenta notablemente el costo, pero no aporta necesariamente más información; la context complexity debería determinarse por la información que contienen los datos, no por un único block size.}

\subsection{Multi-context SAFFU: ramas de parámetros para distintas longitudes}

La sección anterior mostró que una ventana máxima fija hace que los ejemplos cortos paguen quadratic work inútil; la respuesta del multi-context design no es forzar todos los documentos dentro del mismo block, sino que el modelo elija, según la longitud del ejemplo, la rama más pequeña que lo contenga. Al leer el diagrama de arquitectura, compare primero el context width de cada branch y dónde se comparten parámetros; después observe cómo convergen en el mismo espacio de salida. El curso asigna $W_b,U_b$ a distintos context widths $b\in\mathcal B$ y al final los proyecta a una salida compartida. Si la longitud del ejemplo es $K_m$, se elige la rama más pequeña que cumpla $b\ge K_m$, lo que puede escribirse como
\[
b^*(m)=\min\{b\in\mathcal B:b\ge K_m\},
\]
\[
\hat y_m=\operatorname{softmax}\!\left(
\operatorname{softmax}(x_{m,h}W_{b^*}X_m^\top)X_mU_{b^*}O
\right).
\]

\lecturefigure{slide-12-multicontext-design.jpg}{Multi-context SAFFU: elección de rama según el context width}{Grabación oficial de Stanford Online 00:36:29}

\begin{importantbox}{Beneficio y costo de las ramas dinámicas}
El beneficio es que los ejemplos cortos no pagan el quadratic cost completo de la ventana más larga; el costo es mantener varios conjuntos de parámetros o parámetros fraccionables, construir lotes según la longitud y asegurar que las salidas de las distintas ramas caigan en espacios de representación compatibles.
\end{importantbox}

\begin{lstlisting}[caption={Pseudocódigo de agrupación en lotes según la longitud con dynamic context}]
lengths = measure_document_lengths(dataset)
buckets = bucket_by_supported_context(lengths, widths=[128, 256, 512, 1024])

for width, batch in buckets:
    x = pad_only_to(batch, width)
    weights = saffu_parameters.for_width(width)
    loss = language_model_loss(x, weights)
    update_trainable_layers(loss)
\end{lstlisting}

\teachervoice{00:34:55--00:37:34, el ponente explica que lo esencial de la multi-context branch es que la modified attention matrix permite tomar un subconjunto de parámetros según la longitud; la proyección Q/K estándar de baja dimensión no corresponde de forma natural a este tipo de fraccionamiento.}

\subsection{Embedding brittleness: una codificación única no equivale a una representación relacionada}

La unicidad del bit-cipher hace reproducible la inicialización, pero que dos palabras compartan algunos bits no equivale necesariamente a que estén relacionadas semánticamente. El curso propone mejorar el embedding con radial co-occurrence: contar coocurrencias en distintos radios alrededor del token y luego agregar esas correlaciones sobre el deterministic code.

\lecturefigure{slide-13-embedding-brittleness.jpg}{De la bit identity al radial co-occurrence embedding}{Grabación oficial de Stanford Online 00:37:34}

La matriz de coocurrencia del radius $r$ puede denotarse $B^{(r)}$; una forma abstracta de agregación es
\[
E_i=\operatorname{normalize}\!\left(
\operatorname{concat}_{r\in\mathcal R}
\sum_j B^{(r)}_{i,j}e_j
\right).
\]
Aquí $e_j$ es el bit-cipher identity vector, y solo $E_i$ absorbe las relaciones estadísticas locales. La aggregation y el weighting del artículo real deben reproducirse según su implementación; esta fórmula no los sustituye.

\begin{warningbox}{No confunda un «bit interpretable» con un eje semántico}
La regla de encendido de cada bit la determina el algoritmo de codificación y no tiene por qué corresponder a conceptos humanos como noun, tense o sentiment. Si se forma o no un significado debe verificarse con probing, nearest neighbors y tareas posteriores.
\end{warningbox}

\teachervoice{00:37:34--00:39:41, el ponente señala que la naive bit assignment, aunque estable, no elimina la embedding brittleness; se necesitan estadísticas como la radial co-occurrence para volver a introducir la correlación en la representación.}

\subsection{Caching comparisons: la condición esencial de la ganancia de velocidad}

Si la embedding matrix $E$ es fija, las context comparisons de la primera capa
\[
C_{i,j}=e_i^\top e_j
\]
dependen solo de las token identities y pueden precalcularse fuera de línea en una lookup table. Durante el entrenamiento o la inferencia, $C$ se lee directamente a partir de los token IDs y después se aplican $W$ y softmax. Así se ahorran dot products repetidos, pero no el weighting, la aggregation ni la red de las capas superiores.

\lecturefigure{slide-14-caching-vector-comparisons.jpg}{Almacenamiento en caché de las vector comparisons de la primera capa tras congelar el embedding}{Grabación oficial de Stanford Online 00:39:41}

\begin{importantbox}{Cache invalidation rule}
En cuanto el embedding se actualiza a $E'\neq E$, la caché anterior $C=EE^\top$ deja de ser igual a $E'E'^\top$. Por ello, durante el entrenamiento hay que congelar el embedding o reconstruir la cache periódicamente, e incluir el costo de reconstrucción en el tiempo total de entrenamiento.
\end{importantbox}

\begin{knowledgebox}{No se puede omitir la jerarquía de memoria}
La comparison table completa de $V\times V$ puede ser mucho mayor que el propio embedding; no debe suponerse que reside entera en memoria. La implementación requiere caché dispersa o por bloques, gather por lote, o almacenar en caché solo los pares de alta frecuencia. La CPU DRAM (Dynamic Random-Access Memory, memoria principal) tiene más capacidad pero mayor latencia; la GPU HBM (High Bandwidth Memory, memoria de video de alto ancho de banda) es más rápida pero más cara; la SRAM en chip (Static Random-Access Memory, caché/memoria compartida) es la más rápida pero la más pequeña.
\end{knowledgebox}

\teachervoice{00:39:41--00:41:48, Williams describe la caché como el principal ahorro de costo del enfoque near-shallow: cuando el embedding no se actualiza, las comparisons de la capa inferior pueden precalcularse; en cuanto los vectors se descongelan (thaw) y se actualizan, el supuesto de la caché deja de valer.}

\subsection{Packing y dynamic context: una técnica de eficiencia no es un context model}

Packing coloca varios documentos cortos en un mismo block largo y usa una mask para impedir la attention entre documentos. Reduce el padding y el número de lotes, pero sigue construyendo una quadratic matrix de longitud $K_{\max}$. Dynamic context, en cambio, envía los documentos cortos a una branch más corta, con un costo aproximado de
\[
C_{\text{dynamic}}\propto \sum_{m=1}^{B}K_m^2d,
\qquad
C_{\text{fixed}}\propto BK_{\max}^2d.
\]

\lecturefigure{slide-15-packing-contexts.jpg}{Compromisos entre packing, padding y context dinámico}{Grabación oficial de Stanford Online 00:42:14}

\begin{warningbox}{El curso no presenta un matched packing benchmark}
Que dynamic context sea más rápido depende de la kernel efficiency, la bucket fragmentation, el tamaño de lote, la mask implementation y el hardware. En la sesión de Q\&A el ponente dijo explícitamente que aún no lo ha comparado directamente con el oracle packing en modelos grandes; por lo tanto, aquí solo cabe conservar el análisis del mecanismo y el juicio de investigación.
\end{warningbox}

\teachervoice{00:42:38--00:45:55, el ponente considera que packing llena la ventana con documentos no relacionados: es una organización del context eficiente pero no fiel; la dynamic length batching hace que los documentos cortos usen solo un subconjunto más pequeño de parámetros.}

\teachervoice{01:16:20--01:19:00, en la sesión de Q\&A reconoce además que no se realizó una matched large-model comparison; la ganancia relativa de packing frente a dynamic blocks todavía requiere experimentos específicos.}

\subsection{Resumen de la sección}

La eficiencia near-shallow proviene de dos cosas: reducir el quadratic work inútil según la longitud real y reutilizar las comparisons de la capa inferior cuando el embedding es fijo. Ninguna es una optimización gratuita: las ramas dinámicas aumentan la complejidad del modelo y de la formación de lotes, y la caché introduce problemas de capacidad e invalidación.
\section{Precision Language Models: configuraciones, etapas de entrenamiento y velocidad}

La clase llama Precision Language Models a los sistemas que usan SAFFU, inicialización explícita y entrenamiento con pocos datos. Aquí, precision subraya que el punto de partida de los parámetros y los datos de la tarea son «más específicos», no la precisión numérica FP16/BF16. Al leer la tabla de configuraciones hay que considerar al mismo tiempo vocabulary, embedding, context, parameter count, token budget y hardware; no basta con fijarse en el apodo del modelo.

\subsection{Big y potato: dos reference points con objetivos completamente distintos}

La big configuration usa un vocabulary de unos 16K tokens, un context de unos 2K y unos 50M parameters, y se entrena en una antigua GPU Titan V de 12GB; la potato configuration usa un vocabulary de unos 1K tokens, un context corto y unos 150K parameters, y se entrena en tiempo real en una CPU Le Potato con 2GB de RAM. La primera pone a prueba la escala del modelado de lenguaje; la segunda, el scratch learning local.

\lecturefigure{slide-16-plm-configs.jpg}{Reference configurations big y potato de los precision language models}{Grabación oficial de Stanford Online 00:46:08}

\begin{knowledgebox}{PLM no es precision numérica}
En esta clase, PLM es el nombre que el equipo le da a una clase de modelos; no se refiere a mixed precision, FP8 ni quantization. Si un documento solo dice «precision model», debe aclarar que se trata de una etiqueta de architecture/training, para no confundirla con la numerical precision.
\end{knowledgebox}

\begin{warningbox}{El «entrenamiento en tiempo real» debe ir ligado al hardware y a la tarea}
En potato, «tiempo real» se refiere a una binary switch task local, con modelo, vocabulario, context y datos muy pequeños; no puede extrapolarse a entrenar en tiempo real un modelo conversacional de propósito general con 2GB de RAM.
\end{warningbox}

\teachervoice{00:46:14--00:49:04, Williams dice que el equipo aún no sabe cómo se comportaría la arquitectura con trillions of tokens o con RLHF; el objetivo de potato es aprovechar la estabilidad con pocos datos para explorar la capacidad de edge scratch-training, a la que los grandes modelos tradicionales no dan prioridad.}

\subsection{De micro a mega: el verdadero uso de la tabla es auditar recursos}

La tabla de configuraciones presenta varias combinaciones de vocabulary, context, dimensiones embedding/hidden y total parameters. Su valor no está en afirmar que cada fila se entrenó por completo, sino en mostrar cómo crecen los parámetros con el ancho y el context, y en dar nombre a las pruebas de estrés small/medium posteriores.

\lecturefigure{slide-17-saffu-size-table.jpg}{Configuraciones de parámetros y tokens de SAFFU/PLM en distintas escalas}{Grabación oficial de Stanford Online 00:48:22}

Si las matrices principales de una branch son aproximadamente $W\in\mathbb{R}^{d\times d}$, $U\in\mathbb{R}^{d\times h}$ y $O\in\mathbb{R}^{h\times V}$, el total de parámetros puede escribirse de forma aproximada como
\[
P\approx Vd+d^2+dh+hV+P_{\text{other}}.
\]
Conforme crecen el vocabulary $V$ y la dimensión hidden $h$, embedding/output suelen convertirse en los términos dominantes; al aumentar las context branches, $W_b,U_b$ también se replican o se amplían.

\begin{knowledgebox}{Las cuatro cuentas que debe completar la tabla de configuraciones}
\begin{enumerate}
  \item Número de parámetros: ¿qué matrices se comparten y cuáles se replican por context branch?
  \item Tokens de entrenamiento: ¿cuántos datos ve cada etapa (warm start, frozen stage, thawed stage)?
  \item Memoria: ¿cuánto ocupan weights, optimizer state, activations y cache?
  \item Tiempo: ¿se contabilizan la construcción de estadísticas, la backpropagation, el cache rebuild y la evaluation?
\end{enumerate}
Optimizer state (estado del optimizador) se refiere a tensores adicionales como el primer momento $m$ y el segundo momento $v$ de Adam/AdamW, que aumentan de forma considerable la memoria de entrenamiento; no basta con estimarla a partir del tamaño del archivo de parámetros.
\end{knowledgebox}

\subsection{Medium PLM: tres etapas, warm, freeze y thaw}

El medium experiment usa unos 12M parameters y la escala de 100M tokens de BabyLM. Las tres etapas de la slide son: primero, construir un warm start con parte de los datos; después, congelar los vectors y entrenar rápidamente el resto de los parámetros con la mayor parte de los datos; por último, hacer thaw de los vectors y aplicar un fine-tuning conjunto con un learning rate menor.

\lecturefigure{slide-18-medium-plm-training.jpg}{Dinámica de entrenamiento warm-start, frozen-vector y thawed-vector del medium PLM}{Grabación oficial de Stanford Online 00:49:09}

Puede abstraerse como
\[
\theta^{(0)}=\operatorname{ExplicitInit}(D_{\text{warm}}),
\]
\[
\theta^{(1)}=\operatorname{Train}(D_{\text{main}},\theta^{(0)};E\ \text{frozen}),
\]
\[
\theta^{(2)}=\operatorname{FineTune}(D_{\text{tail}},\theta^{(1)};E\ \text{trainable}).
\]

\begin{importantbox}{Qué implican freeze y thaw para el sistema}
Hacer freeze del embedding no solo reduce la memoria de gradient/optimizer, sino que también garantiza que la comparison cache siga siendo válida; el thaw permite que las representaciones se adapten a la tarea final, pero obliga a recalcular las comparisons y puede alterar la feature distribution que corresponde a la solución explícita anterior.
\end{importantbox}

\teachervoice{00:49:09--00:51:35, el ponente explica los 100M tokens de BabyLM como una developmentally plausible exposure, y dice que congelar los vectors mejora a la vez la velocidad de entrenamiento y la estabilidad de la caché, mientras que el thaw final cambia un costo adicional por calidad final.}

\subsection{How fast: una comparación headline debe desglosar sus condiciones}

La slide pone juntos el dato aproximado de unas 100 horas de GPT-2 en 8 A100 y el avance del medium PLM en una antigua Titan V, de unos 18 tokens/hour. Ambos no están estrictamente emparejados en context, vocabulary, implementation, calidad objetivo ni proceso de datos, por lo que solo pueden leerse como motivation y no como un benchmark justo.

\lecturefigure{slide-19-plm-training-speed.jpg}{Referencia aproximada de la clase sobre la velocidad de entrenamiento de PLM}{Grabación oficial de Stanford Online 00:53:21}

Una contabilidad de tiempo más completa debería escribirse como
\[
T_{\text{total}}=T_{\text{stats}}+T_{\text{cache}}
+T_{\text{backprop}}+T_{\text{rebuild}}+T_{\text{eval}},
\]
en lugar de reportar solo el loop principal de entrenamiento. Aunque las estadísticas explícitas y la cache se calculen en la etapa de preprocesamiento de datos, siguen siendo un costo del sistema de entrenamiento.

\begin{warningbox}{No difundir «32 veces más capacidad, por lo tanto 32 veces más velocidad»}
El GPU theoretical throughput, la model FLOPs utilization, el memory bandwidth, el batch size y la kernel maturity modifican el wall-clock. La proporción de hardware de la slide de la clase no sustituye a un profiler ni a una matched implementation.
\end{warningbox}

\teachervoice{00:52:48--00:54:10, el ponente usa GPT-2/A100 como referencia intuitiva; las notas deben conservar su carácter de rough comparison y no reescribirla como un speedup riguroso.}

\subsection{Qué puede hacer un small PLM: distinguir primero fit de transfer}

Al momento de la clase, el equipo había hecho stress tests principalmente hasta cerca de 1B tokens, y muchos experimentos usaron solo de cientos de miles a decenas de millones de tokens. Que un modelo pequeño reduzca rápido la loss de entrenamiento indica que la inicialización es estable; pero la pregunta de fondo es si puede aprender reglas transferibles a partir de application data, y no solo memorizar las muestras locales.

\lecturefigure{slide-20-small-plm-capabilities.jpg}{La pregunta central de investigación del small PLM: ¿puede aprender directamente de los datos de la aplicación?}{Grabación oficial de Stanford Online 00:54:16}

\begin{warningbox}{Small-data fit no equivale a general language ability}
Una PPL muy baja en el conjunto de entrenamiento puede deberse a memorization. Para juzgar si el modelo aprendió una directive semantics robusta se necesitan, como mínimo, pruebas con held-out users, formulaciones no vistas, ASR con ruido, entradas ajenas a la tarea y deriva temporal.
\end{warningbox}

\teachervoice{00:41:57--00:42:33, Williams distingue por iniciativa propia entre «aprender muy rápido unos cuantos miles de tokens» y «convertirse en un modelo de lenguaje capaz de hablar de forma amplia»; lo primero es una observación, lo segundo no queda demostrado por los resultados con pocos datos.}

\teachervoice{00:54:16--00:55:21, el ponente plantea la application-only transfer from scratch como una pregunta abierta, no como una capacidad ya resuelta.}

\subsection{Resumen de la sección}

Las configuraciones PLM van desde un tiny controller hasta un language model mediano, pero cada nivel depende de un vocabulario, un context, unos datos y un hardware distintos. El relato de entrenamiento que de verdad puede reproducirse es warm--freeze--thaw; un headline de velocidad solo tiene sentido de ingeniería si contabiliza por completo statistics, cache y evaluation.
\section{Application-only edge learning: que el interruptor local recopile sus propios datos}

La segunda mitad de la clase aterriza la arquitectura abstracta en una tarea mínima: un interruptor controlado por voz. El objetivo no es construir un assistant de propósito general, sino verificar si el modelo puede recopilar desde cero user data local en un entorno air-gapped, aprender una binary action y completar el entrenamiento y la inferencia en un microprocessor.

\subsection{Sin preentrenamiento general: de dónde salen los datos de la tarea}

El application-only pipeline construye datos supervisados a partir de las interacciones reales del usuario. El sistema escucha la voz, aplica automatic speech recognition (ASR, reconocimiento automático del habla), pide al usuario que accione el interruptor cuando el modelo no está seguro y después empareja la transcript con la action. Así evita reunir primero un corpus de instrucciones a gran escala, pero debe resolver la transcription errónea, el label delay y una distribución de conductas extremadamente estrecha.

\lecturefigure{slide-21-application-only-training.jpg}{Flujo completo para entrenar un PLM solo con datos de interacción de la aplicación}{Grabación oficial de Stanford Online 00:55:21}

\begin{importantbox}{El significado preciso de air-gapped}
Air-gapped significa que el entrenamiento y la inferencia del dispositivo no dependen de una conexión de red y que los datos no salen del entorno local. Eso no equivale automáticamente a seguridad: los permisos del micrófono, las actualizaciones de firmware, el acceso físico, los archivos del modelo y los registros siguen requiriendo un modelado de amenazas.
\end{importantbox}

\teachervoice{00:55:21--00:58:34, el ponente divide la tarea en listen, transcribe, operate y learn: no existe un generic pretraining previo, y la label proviene de que el usuario accione realmente el interruptor cuando hace falta.}

\subsection{Self-supervised interaction: aquí el «self» proviene de la retroalimentación del entorno}

La Slide usa el término self-supervised, pero en sentido estricto el sistema sí recibe action labels; la diferencia es que las etiquetas surgen de forma natural de la interacción local, no de un corpus anotado manualmente de antemano. Si la transcript feature es $x_t$ y la acción del usuario es $y_t\in\{0,1\}$, puede usarse binary cross-entropy
\[
\mathcal L_t=-y_t\log p_\theta(y_t=1\mid x_t)
-(1-y_t)\log\bigl(1-p_\theta(y_t=1\mid x_t)\bigr).
\]

\lecturefigure{slide-22-self-supervised-voice.jpg}{Descomposición listen--transcribe--operate--learn de la tarea de interacción por voz}{Grabación oficial de Stanford Online 00:56:53}

\begin{warningbox}{El ASR error se convierte en label noise}
El usuario pronuncia una oración, pero lo que ve el entrenamiento del modelo es la ASR transcript. Si la transcripción es errónea y la etiqueta de acción es correcta, el par de datos asocia un texto equivocado con la acción. La evaluación debe registrar el word error rate y las condiciones de acento y ruido, y conservar el audio original o información de alineación que permita auditarla.
\end{warningbox}

\subsection{Smart data collection: pedir una acción solo cuando falta información}

La subsección anterior solo enumeró cuatro componentes: listen, transcribe, operate y learn; esta los organiza en un ciclo cerrado que obtiene etiquetas de forma activa. La pregunta central no es «cómo graba el micrófono», sino cuándo confiar en la predicción del modelo, cuándo pedir supervisión al usuario o al interruptor físico y cómo reincorporar las muestras nuevas sin amplificar los errores previos. El sistema no espera pasivamente un dataset fijo, sino que forma durante la interacción un local active-learning loop: al escuchar una instrucción, el modelo emite una predicción; si la confidence es baja o el usuario la corrige, solicita u observa la acción real y la agrega al conjunto de entrenamiento.

\lecturefigure{slide-23-smart-data-collection.jpg}{Ciclo cerrado de aprendizaje local formado por usuario, micrófono, ASR, modelo e interruptor}{Grabación oficial de Stanford Online 00:58:34}

\begin{lstlisting}[caption={Máquina de estados del aprendizaje interactivo del interruptor de voz local}]
while device_is_on:
    audio = microphone.listen()
    transcript = asr.transcribe(audio)
    action, confidence = model.predict(transcript)

    if confidence < threshold or user_overrides(action):
        label = observe_physical_switch()
        local_dataset.append(transcript, label)
        model.update(local_dataset.recent_examples())
    else:
        controller.execute(action)
\end{lstlisting}

\begin{knowledgebox}{Tres riesgos del online learning}
\begin{enumerate}
  \item \textbf{Feedback loop:} los errores del modelo pueden influir en la conducta del usuario y luego tratarse como datos nuevos.
  \item \textbf{Catastrophic forgetting:} entrenar solo con comandos recientes puede hacer que se olviden expresiones anteriores.
  \item \textbf{User drift:} la forma de hablar, el ruido ambiental y el uso del dispositivo cambian con el tiempo.
\end{enumerate}
La clase no ofreció una solución completa a estos problemas; por ello, al reproducir el experimento hay que diseñar además replay, holdout y rollback.
\end{knowledgebox}

\teachervoice{00:58:34--01:00:10, Williams describe la recopilación de datos como un ciclo cerrado de interacción: el sistema escucha y transcribe localmente, y solo pide al usuario que accione el interruptor cuando necesita supervisión, con lo que acumula gradualmente directive/action pairs.}

\subsection{Dialogue representation: escribir la binary action como una secuencia de lenguaje}

El ciclo cerrado de interacción resuelve de dónde vienen las etiquetas; falta decidir en qué forma entran los datos al modelo. La clase no diseñó para el interruptor un classifier completamente distinto, sino que encadenó los comandos del usuario y las acciones del dispositivo en un dialogue, de modo que el mismo character/token language model lea la utterance y prediga un action token estructurado. Al leer la Slide 24 conviene ubicar primero la frontera entre entrada y salida, y después juzgar si la extensibilidad que aporta este formato unificado compensa el costo adicional de la secuencia. Esta representación también puede desperdiciar capacidad en el texto fijo de la plantilla.

\lecturefigure{slide-24-dialogue-data.jpg}{La interacción con el interruptor representada como un user/device dialogue}{Grabación oficial de Stanford Online 01:02:15}

Puede abstraerse como la secuencia
\[
s_t=[\texttt{USER}:u_t,\ \texttt{DEVICE}:a_t],
\]
donde $u_t$ es la transcript y $a_t\in\{\texttt{ON},\texttt{OFF},\texttt{NOOP}\}$. Si el único objetivo es el control, un classifier directo es más ligero; solo si en el futuro se quiere generar una clarification o una reply, la dialogue formulation ofrece margen de extensión.

\begin{warningbox}{Un binary switch no equivale a general dialogue}
El modelo solo debe predecir ON/OFF/NOOP en un action space sumamente estrecho. La diversidad en la formulación de los comandos puede ser alta, pero la complejidad de la salida sigue siendo muy inferior a la de la generación de lenguaje abierta.
\end{warningbox}

\subsection{Potato y orange: la configuración edge aún requiere la contabilidad del sistema completo}

La clase presentó dos configuraciones pequeñas. Potato usa vocabulary, embedding y context más pequeños y se entrena en Le Potato; orange es un poco mayor y busca reducir la latency y mejorar la predicción. El número de parámetros de la tabla es solo el weight count; hay que sumar la memoria de ASR, audio buffering, runtime y la actualización del modelo.

\lecturefigure{slide-25-potato-configs.jpg}{Edge configuration de los PLM potato y orange}{Grabación oficial de Stanford Online 01:02:38}

\begin{knowledgebox}{Latency de extremo a extremo del edge controller}
\[
T_{\text{interaction}}=T_{\text{audio}}+T_{\text{ASR}}+T_{\text{model}}
+T_{\text{controller}}+T_{\text{actuator}}.
\]
Optimizar solo $T_{\text{model}}$ no necesariamente hace que el usuario perciba mayor rapidez. Los varios segundos de retraso de la demostración en clase pueden deberse sobre todo a la transcription, a la CPU inference o a llamadas en serie; hace falta un profiler para desglosarlos.
\end{knowledgebox}

\teachervoice{01:00:10--01:02:48, el ponente explica que el potato PLM character-level más pequeño aprende las directives con unos 20 minutos de interacción; una configuración mayor mejora la predicción, pero aumenta el costo de recursos y de respuesta.}

\subsection{Emergent positive performance: la predicción útil aparece solo cuando hay suficientes muestras}

La tabla de resultados muestra que, al principio, cuando faltan datos, las positive/negative predictions casi carecen de sentido; tras acumular cierto volumen de interacciones, la action prediction correcta sube notablemente. Aquí la emergence es un sample-threshold phenomenon de una tarea pequeña, no la emergent ability entre tareas de los benchmarks de modelos grandes.

\lecturefigure{slide-26-emergent-positive-performance.jpg}{Capacidad de predicción del interruptor que aparece al aumentar las muestras de interacción local}{Grabación oficial de Stanford Online 01:02:48}

\begin{importantbox}{Cómo verificar que no es un artefacto del class imbalance}
Reporte el número de muestras de cada clase ON/OFF/NOOP, la balanced accuracy, precision/recall, la confusion matrix, un holdout separado por usuario y never-seen paraphrases. Si la mayor parte del tiempo el interruptor no cambia, la accuracy por sí sola puede quedar inflada por la majority class.
\end{importantbox}

\teachervoice{01:02:48--01:04:08, Williams dice que al principio, sin datos suficientes, el modelo casi no puede predecir nada; conforme se acumulan muestras, la capacidad da un salto. Sin embargo, en Le Potato cada interacción sigue tardando varios segundos: el prototype es viable, pero obliga a esperar.}

\subsection{Where next: tras la viabilidad queda un largo camino de ingeniería}

La última slide didáctica presenta el PLM local como un edge-trainable controller. El trabajo futuro incluye que el dispositivo haga talk-back, phone actions más generales, un hybrid entre la attention tradicional y SAFFU, un bit-cipher mejorado, warm starts de distintas layers, una larger evaluation y el BabyLM challenge.

\lecturefigure{slide-27-where-next.jpg}{Agenda abierta para pasar del binary switch a un controlador local más general}{Grabación oficial de Stanford Online 01:04:12}

\begin{warningbox}{Al terminar la clase aún no había una implementación pública}
En la sesión de Q\&A se dijo que el código se publicará cuando estén terminados el artículo y la evaluation infrastructure. Ninguna reproducción puede suponer que existe un SDK maduro, un benchmark adapter estándar o un checkpoint descargable directamente.
\end{warningbox}

\teachervoice{01:04:11--01:06:50, el ponente llama al binary switch un punto de partida, no un punto de llegada; el conversation role, las phone actions, la hybrid attention, un better bit-cipher, los larger models y la evaluación estándar siguen siendo work in progress.}

\teachervoice{01:19:02--01:19:41, la sesión final de Q\&A aclara que en ese momento aún no existía una SAFFU implementation pública; el principal obstáculo es que todavía no se han construido las evaluation functions que requiere una architecture no estándar.}

\subsection{Resumen de la sección}

El experimento con Le Potato muestra que una tarea mínima puede arrancar localmente a partir de datos de interacción y producir una prediction útil en unos 20 minutos. Al mismo tiempo, expone limitaciones reales: ASR noise, un action space estrecho, class balance, multi-second latency, la falta de una evaluation estándar y la ausencia de una implementación pública.
\section{Ampliación de Q\&A: alcance, supuestos matemáticos y comparaciones pendientes}

La sesión de Q\&A no agregó ninguna slide propia, pero completó los cuatro límites del texto principal que más fácilmente se exageran: la architecture de la warm-start evidence, la cobertura de las layers que no son softmax, la non-negative assumption y la ausencia de la packing comparison.

\subsection{Las mixed modalities son concebibles, pero cada tipo de layer debe cumplir las condiciones}

Si las image features y las text features pasan por un procesamiento que las vuelve no negativas y entran en una softmax layer, la misma forma de generalized co-occurrence puede, matemáticamente, acumular mixed inputs. Puede escribirse como
\[
H_m=\operatorname{concat}(H_m^{\text{image}},H_m^{\text{text}}),
\qquad F=H^\top Y.
\]
Sin embargo, un multimodal model real también contiene componentes como convolution/vision encoder, normalization, projection y cross-attention; no basta con inicializar la última capa para afirmar que el modelo completo se optimizó explícitamente.

\teachervoice{01:14:33--01:15:25, el ponente considera que las mixed non-negative features pueden compartir el mismo softmax warm start, por lo que image captioning es posible en lo conceptual; esto es solo una dirección que puede derivarse, no un sistema terminado.}

\subsection{Por qué la non-negativity no es un detalle de implementación}

La fórmula contiene $\log F_{j,i}$. Si una feature tiene valores negativos, las outer-product sums pueden ser negativas, y el log deja de corresponder a una probabilidad o a una intensidad de co-ocurrencia. Esto puede tratarse mediante feature splitting, una positive transform o una derivación independiente, pero cada opción cambia el significado estadístico y las condiciones de optimalidad.

\begin{importantbox}{Verificación directa del supuesto de no negatividad}
Antes de implementar la inicialización explícita, revisa primero feature minimum, negative fraction, zero columns y column mass. No tomes en silencio el absolute value cuando aparezcan negativos: eso altera la información de dirección y deja de corresponder a la derivación original.
\end{importantbox}

\teachervoice{01:15:29--01:16:14, Williams explica de dónde viene el supuesto con la pregunta «¿qué ocurre cuando un negative value entra en el logarithm?», y afirma explícitamente que las features negativas requieren una nueva derivation, no un pequeño parche.}

\subsection{Escenarios de uso de la attention estándar, SAFFU y la concatenation}

Los tres mecanismos resuelven problemas distintos. La concatenation conserva un feature slot independiente para cada posición y casi no requiere reponderación, pero su dimensión crece linealmente con el context; la standard attention superpone las distintas posiciones y luego aprende a seleccionar; SAFFU transforma directamente la comparison matrix y realiza la inicialización explícita y el almacenamiento en caché en torno a esa matriz.

\begin{knowledgebox}{Comparación de tres representaciones del context}
\begin{tabularx}{\textwidth}{L{0.19\textwidth} L{0.25\textwidth} X}
\toprule
Método & Ventajas & Principal costo/limitación \\
\midrule
Concatenation & Información posicional completamente separada; fácil de usar en experimentos de una sola capa & La dimensión crece con el context; no es adecuada para secuencias largas \\
Standard attention & Projection flexible, múltiples cabezas, se adapta a distintas representaciones & Muchos parámetros aleatorios, quadratic comparisons \\
SAFFU proposal & Admite inicialización explícita; permite almacenar en caché comparaciones fijas & Supuestos fuertes, implementación no estándar, la composition no está del todo resuelta \\
\bottomrule
\end{tabularx}
\end{knowledgebox}

\subsection{Resumen de la sección}

La sesión de Q\&A devuelve el curso de la «asombrosa historia de eficiencia» a límites verificables: no todos los experimentos son de attention, abarcar otros layer types requiere nuevas derivaciones, las negative features invalidan la fórmula log actual, la comparación de packing aún no se completa y el software sigue en construcción.
\section{Síntesis y ampliación}

Lo que más vale la pena conservar de esta clase no es un speedup headline en particular, sino una forma de pensar que conecta architecture, statistics y systems: si de las estadísticas de los datos pueden obtenerse valores iniciales interpretables, es posible reducir el entrenamiento aleatorio; si la representación subyacente es fija, es posible almacenar en caché las comparaciones repetidas; si la tarea es muy estrecha y los datos se generan localmente, es posible trasladar el entrenamiento al edge. Cada paso trae nuevas restricciones: no se puede tomar solo el beneficio y descartar las premisas.

\subsection{Tres niveles de compresión: inicialización, reutilización del cómputo y ciclo de datos local}

El primer nivel es la initialization: bit-cipher proporciona una deterministic non-negative representation, y la generalized co-occurrence proporciona un warm start para la softmax layer. El segundo nivel es la computation reuse: los frozen embeddings permiten almacenar en caché las pairwise comparisons, y el dynamic context evita que las muestras cortas paguen el costo de la ventana más larga. El tercer nivel es el data loop: el edge controller recopila transcript/action pairs a partir de interacciones reales y cambia una tarea estrecha por la posibilidad de aprender con pocos datos.

\begin{importantbox}{La conclusión mínima confiable del curso}
\begin{enumerate}
  \item Para una softmax layer que cumple los supuestos, un data-derived warm start puede superar a un punto de partida aleatorio.
  \item El iterative training posterior al warm start sigue siendo valioso, sobre todo para la composition de varias capas y para la adaptación de la representación final.
  \item Las fixed representations permiten almacenar en caché parte de las comparaciones, pero la caché, la memoria y la invalidation deben contarse como costos del sistema.
  \item Un tiny local controller puede aprender acciones estrechas a partir de las interacciones con la aplicación, pero de ello no se puede deducir una capacidad de diálogo general.
\end{enumerate}
\end{importantbox}

\subsection{Checklist para reproducir los experimentos}

Al reproducir este tipo de trabajos, lo que más fácilmente se omite son los supuestos y los costos adicionales. La siguiente checklist debe publicarse junto con las gráficas de resultados.

\begin{lstlisting}[caption={Lista de reproducción y publicación de SAFFU/PLM}]
record_official_architecture_and_parameter_shapes()
verify_nonnegative_features_and_softmax_activation()
publish_explicit_initialization_formula_and_epsilon()
report_priming_number_distribution_and_sensitivity()
separate_stats_cache_backprop_and_evaluation_time()
compare_tuned_cold_warm_frozen_and_thawed_baselines()
document_cache_size_location_hit_rate_and_invalidation()
evaluate_heldout_users_paraphrases_noise_and_class_balance()
report_end_to_end_edge_latency_not_model_latency_only()
publish_code_data_splits_seeds_and_failure_cases()
\end{lstlisting}

\begin{warningbox}{Al ver «sin pretraining/backprop», pregunte de inmediato}
¿Qué capa es la que no lo necesita? ¿Cuántos datos usaron las estadísticas explícitas? ¿Se hace gradient tuning después? ¿El embedding está congelado? ¿La prueba mide training fit o held-out generalization? ¿Las unsupported layers se inicializan aleatoriamente? Si estas preguntas no tienen respuesta, el headline es mucho mayor que la evidencia.
\end{warningbox}

\subsection{Preguntas de autoevaluación}

\begin{multicols}{2}
\small
\begin{enumerate}
  \item ¿Qué relación hay, en cuanto a parametrización, entre la attention estándar $QK^\top$ y $XWX^\top$?
  \item ¿Por qué SAFFU sigue siendo una attention con parámetros y no una similitud sin parámetros?
  \item ¿Near-shallow describe la depth, o la actualización en línea y la dependencia de cómputo?
  \item ¿Por qué bit-cipher requiere $b\ge\lceil\log_2(V+1)\rceil$?
  \item ¿En qué se diferencia la generalized co-occurrence $H^\top Y$ de la coocurrencia tradicional en ventanas de palabras?
  \item ¿Por qué la softmax solution explícita requiere non-negative features?
  \item ¿Qué corrige el priming number $K$ en la fórmula?
  \item ¿Por qué los hidden targets de la attention son más difíciles de obtener que los decoder targets?
  \item ¿Por qué la gráfica de PPL con warm start no demuestra directamente que la attention de SAFFU sea mejor?
  \item ¿Qué pueden demostrar los resultados en MNIST y qué no pueden demostrar?
  \item ¿Por qué un frozen embedding afecta a la vez la velocidad de entrenamiento y la cache correctness?
  \item ¿Qué tipo de desperdicio reduce el packing y cuál el dynamic context?
  \item ¿Por qué la precisión de PLM no es FP16/BF16?
  \item ¿Qué optimiza cada una de las tres etapas warm--freeze--thaw?
  \item ¿Qué componentes incluye la latency de extremo a extremo de Le Potato?
  \item ¿Por qué la binary switch accuracy puede ser engañosa debido al class imbalance?
  \item ¿Qué conclusiones seguían siendo solo future work al terminar la clase?
\end{enumerate}
\end{multicols}

\subsection{Lecturas adicionales}

\begingroup
\small
\setlength{\itemsep}{2pt}
\begin{itemize}
  \item Williams and Zhao, \href{https://arxiv.org/abs/2311.07510}{\emph{Explicit Foundation Model Optimization with Self-Attentive Feed-Forward Neural Units}}: artículo principal sobre la arquitectura SAFFU, la inicialización explícita y la ablation.
  \item Williams and Zhao, \href{https://arxiv.org/abs/2311.07498}{\emph{Reducing the Need for Backpropagation and Discovering Better Optima With Explicit Optimizations of Neural Networks}}: solución explícita de softmax de una sola capa, MNIST y local multi-layer warm start.
  \item Zhao and Williams, \href{https://arxiv.org/abs/2311.11012}{\emph{Bit Cipher}}: deterministic bit vectors, co-occurrence aggregation y transfer experiments.
  \item Williams and Heidenreich, \href{https://arxiv.org/abs/2205.00148}{\emph{To Know by the Company Words Keep and What Else Lies in the Vicinity}}: Word2Vec softmax, co-occurrence factorization y la motivación del radial context.
  \item Frankle and Carbin, \href{https://arxiv.org/abs/1803.03635}{\emph{The Lottery Ticket Hypothesis}}: antecedente original para entender la motivación de esta clase de «reducir la lotería aleatoria»; no debe tomarse como conclusión experimental de SAFFU.
  \item Warstadt et al., \href{https://arxiv.org/abs/2310.00915}{\emph{Findings of the BabyLM Challenge}}: para revisar los requisitos concretos de developmentally plausible data scale y de la evaluation estándar.
\end{itemize}
\endgroup

\end{document}
